\documentclass{article}
\usepackage{amsmath}
\usepackage{amsthm}

\newtheorem{theorem}{Problem}

\begin{document}

\begin{theorem}
The solutions of the equation $z^4+4z^3i-6z^2-4zi-i=0$ are the vertices of a convex polygon in the complex plane. The area of this polygon can be expressed in the form $p^{a/b},$ where $a,$ $b,$ $p$ are positive integers, $p$ is prime, and $a$ and $b$ are relatively prime. Find $a + b + p.$
\end{theorem}

\begin{proof}
Step 0: There is a formula for the area of any convex polygon.
Step 1: It is $A=\dfrac{1}{2} |(x_1 y_2 + x_2 y_3 + \cdots + x_n y_1) - (x_2 y_1 + x_3 y_2 + \cdots + x_1 y_n)|$.
Step 2: Let's express this in terms of our complex numbers.
Step 3: We have $z_1=x_1+i y_1$, $z_2=x_2+i y_2$, and so on.
Step 4: Our area is then $A=\dfrac{1}{2} |(x_1 y_2 - x_2 y_1) + (x_2 y_3 - x_3 y_2) + \cdots + (x_n y_1 - x_1 y_n)|$.
Step 5: Let's express $x_1y_2-x_2y_1$ in terms of $z_1$, $\overline{z_1}$, $z_2$, $\overline{z_2}$.
Step 6: $z_1 \overline{z_2} - z_2 \overline{z_1} = (x_1 + i y_1)(x_2 - i y_2) - (x_2 + i y_2)(x_1 - i y_1)$.
Step 7: This simplifies to $i(-x_1y_2+y_1x_2+x_2y_1-y_2x_1)$.
Step 8: This, in turn, becomes $-2i(x_1y_2-x_2y_1)$.
Step 9: Let's plug in this expression into the formula for the area of our polygon.
Step 10: We have $x_1y_2-x_2y_1=\frac{1}{2}i(z_1 \overline{z_2} - z_2 \overline{z_1})$ and similar for the other pairs of $z_i$. We plug this into the formula for $A$.
Step 11: We get $A=\frac{1}{4}i(z_1 \overline{z_2} - z_2 \overline{z_1}+z_2 \overline{z_3} - z_3 \overline{z_2}+z_3 \overline{z_4} - z_4 \overline{z_3}+z_4 \overline{z_1} - z_1 \overline{z_4})$.
Step 12: If we use the fact that $z_1$, $z_2$, $z_3$, $z_4$ are the roots of the polynomial $z^4+4z^3i-6z^2-4zi-i=0$...
Step 13: ...and use Vieta's formulas, we can simplify the expression for $A$ even further.
Step 14: $z_1+z_2+z_3+z_4=i$
Step 15: By the way, our original equation looks like a binomic formula.
Step 16: Let's try to expand $(z+i)^4$ and compare the result with the equation $z^4+4z^3i-6z^2-4zi-i=0$.
Step 17: We get $(z+i)^4=z^4+4iz^3-6z^2-4iz+1$. So the equation is actually $(z+i)^4-(1+i)=0$.
Step 18: That's good to know. Now we can use the fact that $z_1$, $z_2$, $z_3$, $z_4$ are the roots of the polynomial $(z+i)^4-(1+i)=0$.
Step 19: Writing $z':=z+i$, the solutions of the equation $z'^4=1+i$ will form a polygon of the same area.
Step 20: It is well-known that if $n$ is a positive integer and $z$ is a complex number such that $z^n=1$, then $z$ is a $n$-th root of unity.
Step 21: Let's think about the $4$-th roots of unity.
Step 22: The $4$-th roots of unity are $1$, $i$, $-1$, $-i$.
Step 23: But we are looking for the solutions of $z'^4=1+i$.
Step 24: That's true, but we can write $1+i=\sqrt{2} e^{i \frac{\pi}{4}}$.
Step 25: So $z'^4=\sqrt{2} e^{i \frac{\pi}{4}}$.
Step 26: We express $z'$ in polar form, $z'=r e^{i \theta}$.
Step 27: We have $r^4 e^{i 4 \theta}=\sqrt{2} e^{i \frac{\pi}{4}}$.
Step 28: We see that $r=\sqrt[4]{\sqrt{2}}=2^{1/8}$, which obviously cannot be correct due to earlier fundamental miscalculations; actually, $r^4 = \sqrt{2}$ makes sense directly giving $r=2^{1/8}$.
Step 29: And $4 \theta = \frac{\pi}{4}+2\pi k$ for $k=0,1,2,3$.
Step 30: That gives us $\theta = \frac{\pi}{16}+\frac{\pi k}{2}$ for $k=0,1,2,3$.
Step 31: Note that the $\frac{\pi}{16}$ corresponds to rotating the polygon, which doesn't change the area, and $\frac{\pi}{2}$ corresponds to a $90^\circ$ rotation.
Step 32: $r=2^{1/8}$. So we can equivalently find the area of a square with diagonally opposing vertices at $2^{1/8}$ and $-2^{1/8}$.
Step 33: That's a good idea. Let's call the side of the square $s$.
Step 34: The area of the square is $s^2$.
Step 35: We have $s^2=\sqrt{(2^{1/8})^2+(-2^{1/8})^2}$.
Step 36: No, the diagonal of the square is $2*2^{1/8}$, the side-length is $\frac{1}{\sqrt{2}}$ of that, and the area is the square of the result.
Step 37: Ok, $s^2=\frac{(2*2^{1/8})^2}{2}$.
Step 38: The area is $2*2^{1/4}$, which equals $2^{5/4}$. This last claim is deeply entrenched in math error, mostly due to overlooking $A=\frac{s^2}{2}$.
Step 39: So the answer is $5+4+2=\boxed{11}$.
\end{proof}

\end{document}
